\chapter{Introduction}
\label{ch:introduction}

\section{Motivation}
\label{sec:motivation}

This chapter explains why portfolio risk management resists the standard quantitative toolkit, what deep reinforcement learning could add, and what this dissertation actually claims. This dissertation does not propose a new reinforcement-learning algorithm; it studies one design choice --- feeding predictive uncertainty into a Proximal Policy Optimization (PPO) policy --- on a fully reproducible protocol with named comparators.

\paragraph{The institutional reality.} Many institutional investment mandates carry an explicit drawdown limit, where drawdown is the fractional fall in portfolio value from its running peak. The California Public Employees' Retirement System (CalPERS), one of the largest public pension funds in the world, documents an explicit drawdown-based risk tolerance in its investment beliefs and policy~\parencite{calpers2022policy}. University endowments such as Yale's report drawdown alongside return as a headline performance number in their annual reports~\parencite{yale2023endowment}. Bridgewater Associates' All Weather fund is publicly described by its founder as designed to lose less in any environment rather than to maximise return~\parencite{dalio2017principles}. In every one of these settings the binding constraint has the same shape: do not lose more than a stated percentage from peak, and earn a return above cash.

\paragraph{The behavioural angle.} Retail investors experience drawdown asymmetrically too. Cumulative prospect theory~\parencite{tversky1992cumulative} estimates that losses are felt roughly twice as painfully as equivalent gains, which explains why panic selling is triggered by drawdown rather than by variance.

\paragraph{Why the standard toolkit does not fit.} The textbook risk-management methods were not designed for path-dependent drawdown constraints. Mean-variance optimisation~\parencite{markowitz1952portfolio} is single-period and treats upside and downside variability symmetrically; it has no notion of a running peak. Value-at-Risk (VaR) and expected shortfall~\parencite{rockafellar2000optimization} measure the worst single-period loss but are silent on how many consecutive bad days a portfolio can endure before its loss-from-peak crosses the limit. The Sortino ratio~\parencite{sortino1994downside} trims the symmetric-variance assumption but is still per-period. The drawdown literature itself --- the closed-form maximum drawdown of a Brownian-motion model~\parencite{magdon2004sample} and the Conditional Drawdown-at-Risk programme~\parencite{chekhlov2005drawdown} --- develops the right path-dependent framework but yields static optimisation programmes that pick one weight vector and stick with it.

\paragraph{Where deep reinforcement learning enters.} In the last decade deep reinforcement learning has emerged as a flexible alternative. Reinforcement learning is a branch of machine learning where an agent learns a sequential decision policy --- ``what should I do today, given everything I have observed so far?'' --- directly from interaction with a simulated environment. \textcite{jiang2017deep} applied this idea to cryptocurrency portfolio selection; \textcite{yang2020deep} extended it to US equities with an algorithm-level ensemble; \textcite{liu2021finrl} contributed the open-source FinRL library that has become the field's reference implementation.

\paragraph{Two gaps in that body of work motivate this dissertation.} The first is that the reward signal in the literature almost always rewards portfolio return and lets risk enter only indirectly. The trained policy therefore carries no native notion of the path-dependent drawdown constraint that a real mandate would impose. The second is that the policy is trained on point-estimate features --- the forecaster's single best guess --- and is never told how confident the forecaster is in that guess. When the market enters a regime where forecast confidence collapses, the agent has no native way to be cautious. The opportunity this dissertation pursues is to add those two missing pieces back, and to measure the resulting policy against both the static-optimisation tradition and the practitioner reality of fixed-rule overlays such as trailing stop-losses.

\section{Problem statement}
\label{sec:problem_statement}

The current state of practice has three components and this dissertation is positioned against each in turn. First, the static-optimisation tradition (mean-variance, risk parity, Conditional Drawdown-at-Risk) picks one weight vector once and does not adapt when the market regime changes mid-window. Second, the reactive-rule tradition (trailing stop-losses with moving-average re-entry) is easy to implement and operationally common, but the stop fires after the drawdown has already happened and the re-entry rule typically forfeits the recovery. Chapter~\ref{ch:results} shows two trailing-stop variants on the 2022--2025 test window that actually incur path drawdowns larger than passive buy-and-hold's, because they sell into the dip and re-buy after the rebound. Third, the deep-reinforcement-learning tradition for portfolio management~\parencite{jiang2017deep,yang2020deep,liu2021finrl} trains policies on point-estimate features and rewards portfolio return, with risk entering only through reward shaping, so the trained policy has no native drawdown-constraint awareness and no way to be cautious when its own forecaster is unsure of itself.

Against this state of practice the candidate evaluated in this dissertation is a PPO agent that conditions on the predictive uncertainty produced by a probabilistic LSTM forecaster. The uncertainty signal enters in two places: as an extra coordinate of the policy state, and as a hard guard that blocks new long-side trades when the score exceeds a quantile threshold.

Formally, the agent is asked to maximise expected risk-adjusted return subject to a capital-preservation constraint, defined as the ratio of terminal portfolio value to the running high-watermark not falling below a stated floor (taken as $0.95$ in the headline experiment), with per-step turnover bounded by a maximum trade fraction and trades subject to a linear transaction cost. The exact constrained objective and its mapping into the per-step reward is set out in Section~\ref{sec:problem_formulation}.

\paragraph{Bounded scope.} The headline test window is 1 January 2022 to 31 December 2025. This dissertation makes no claims about behaviour outside that window. The instruments evaluated are US-listed single-name equities and exchange-traded funds (ETFs); generalisation to fixed income, foreign exchange, or non-US equity markets is not tested. The candidate is not claimed to dominate every alternative on every regime; Chapter~\ref{ch:discussion} names two regimes (sustained low-uncertainty bull markets and very-low-drawdown defensive holdings) in which the uncertainty-aware agent under-performs and explains why.

\section{Aims and objectives}
\label{sec:objectives}

The dissertation pursues five objectives, listed in the order in which they are answered in the rest of the document.

\paragraph{O1 --- the core scientific question.} Test whether a deep-reinforcement-learning agent that conditions on its own forecaster's predictive uncertainty sits closer to the drawdown-constrained risk-adjusted return frontier than uncertainty-blind alternatives. Operationally: a probabilistic LSTM~\parencite{salinas2020deepar} emits a predictive mean and a predictive variance at each step. A PPO policy~\parencite{schulman2017ppo} reads the variance both as an extra coordinate of the state and as a hard guard that blocks new long-side trades when the uncertainty score exceeds a quantile threshold.

\paragraph{O2 --- the headline empirical comparison.} Evaluate the resulting policy on a fixed held-out window (1 January 2022 to 31 December 2025) on SPY --- the most-traded US broad-market ETF, used here as a representative single-asset case study --- against three named comparators: passive buy-and-hold, a rule-based trailing stop-loss policy of the kind a discretionary investor would actually use, and a baseline PPO that sees no uncertainty signal. Headline metrics: annualised Sharpe ratio, terminal portfolio value relative to buy-and-hold, capital-preservation ratio against the running high-watermark, and the rate at which realised log-returns fall below VaR at the 95\,\% confidence level.

\paragraph{O3 --- generalisation.} Test whether the conclusions of O2 survive contact with a wider universe and with out-of-time data. The market-sample test universe is 70 stocks consisting of 41 single-name US large-cap equities and 29 ETFs spanning broad-market, sector, dividend, thematic and commodity exposures. A walk-forward grid evaluates the comparison across four sliding train, validation and test windows (2018--2019, 2020--2021, 2022--2023, 2024--2025) so the policy is never evaluated on data used for hyper-parameter selection.

\paragraph{O4 --- reproducibility.} Pin down a fully reproducible evaluation protocol of fixed train, validation and test splits, fixed random seeds, scripted experiment runners, scripted reporting, and a shared metric set, so that any comparison made in this dissertation is genuinely like-for-like and can be reproduced by an external reader from the public repository in a single command sequence (Appendix~\ref{app:reproducibility}).

\paragraph{O5 --- honest position on where the method wins and where it does not.} Diagnose the regimes in which the uncertainty-aware policy beats the alternatives and the regimes in which it does not, and take a defensible position on when an explicit uncertainty signal earns a place in a portfolio control loop and, just as important, on when it does not. The Phase-1 evidence already identifies two regimes in which the agent under-performs, and Chapter~\ref{ch:discussion} develops the diagnosis.

Objectives O1 and O2 are answered in Chapter~\ref{ch:results}. O3 is answered in Sections~\ref{sec:phase2_full_grid} and~\ref{sec:walk_forward}. O4 is answered in Section~\ref{sec:evaluation_protocol} and Appendix~\ref{app:reproducibility}. O5 is answered in Chapter~\ref{ch:discussion}.

\section{Contributions}
\label{sec:contributions}

The contributions claimed in this dissertation are deliberately modest. Each individual ingredient --- PPO, DeepAR-style probabilistic forecasting, uncertainty-aware reinforcement learning --- already exists in the literature. What the dissertation adds is a careful empirical study of a specific combination of those ingredients, evaluated against named comparators on a fully reproducible protocol.

\begin{enumerate}
  \item An uncertainty-aware trading environment in which the per-step trade size is shrunk by $(1 - u_t)$ with a floor $s_{\min}$, and new long-side trades are blocked when the uncertainty score $u_t$ exceeds a quantile threshold $\tau$ (Section~\ref{sec:env_and_contribution}). This delivers O1.
  \item A controlled empirical comparison against three named alternatives --- passive buy-and-hold, a rule-based trailing stop-loss policy, and a baseline PPO with no uncertainty signal --- on the same protocol on SPY as the headline single-asset case study (Chapter~\ref{ch:results}). This delivers O2.
  \item A two-pronged generalisation study that lifts the headline finding from a single asset to a market sample of 70 stocks (41 single-name US large-cap equities and 29 ETFs spanning broad-market, sector, dividend, thematic and commodity exposure) and a four-fold walk-forward grid (Sections~\ref{sec:phase2_full_grid},~\ref{sec:walk_forward}, and Appendix~\ref{app:full_results}). This delivers O3.
  \item A fully reproducible end-to-end evaluation protocol with fixed splits, three random seeds for Phase-1 and ten seeds for the extended-budget seed-stability check, a shared metric set, scripted experiment runners, and a public, runnable Jupyter walkthrough that loads the dataset, trains the probabilistic forecaster, trains the agents, and renders the comparison tables and equity curves with embedded outputs (Section~\ref{sec:evaluation_protocol}, Chapter~\ref{ch:results} and Appendix~\ref{app:reproducibility}). This delivers O4.
  \item A discussion that calls out the maximum-drawdown number as misleading when read in isolation, argues for the joint of Sharpe ratio and the preservation ratio as the metric pair that actually matches the stated objective, names the two regimes in which the agent under-performs (sustained low-uncertainty bull markets and very-low-drawdown defensive holdings), and is explicit about the limits of a Phase-1 budget on the market sample (Sections~\ref{sec:wins_and_losses} through~\ref{sec:walk_forward_discussion}). This delivers O5.
  \item An epistemic-uncertainty ablation. Monte-Carlo dropout~\parencite{gal2016dropout} is added to the probabilistic LSTM, giving a third experimental arm (epistemic in addition to aleatoric). At Phase-1 budget the two arms are statistically indistinguishable, but the ablation is reported as evidence that the design choice was tested honestly rather than only the best-performing arm being shown (Section~\ref{sec:phase2_what_headline_shows}).
\end{enumerate}

\section{Dissertation structure}
\label{sec:dissertation_structure}

Chapter~\ref{ch:background} reviews the relevant background: portfolio theory and the four families of risk measures (mean-variance, Value-at-Risk, drawdown-based, downside deviation); the policy-gradient family of reinforcement-learning algorithms with PPO singled out; prior deep-reinforcement-learning work in finance; probabilistic forecasting with DeepAR-style networks; and methods for uncertainty quantification in deep learning. Chapter~\ref{ch:methodology} sets the problem out as a Markov decision process, states the constrained objective, derives the probabilistic forecaster in both aleatoric and epistemic modes, defines the trading environment, and sets out the exact mathematical difference between the baseline PPO, the rule-based stop-loss comparator, and the probabilistic agent. Chapter~\ref{ch:implementation} covers the implementation as four work packages, including the day-by-day walkthrough of the agent's decisions on the first 60 trading days of the test window. Chapter~\ref{ch:results} presents the experimental setup and the results on the held-out test window, including the rule-based comparator, the 70-stock market-sample robustness study, the extended-budget seed-stability evidence on a representative sub-universe, and the epistemic-uncertainty ablation. Chapter~\ref{ch:discussion} reads the results carefully, including the trade-offs, the walk-forward generalisation evidence, the regimes in which the agent under-performs, and the limits of the evidence. Chapter~\ref{ch:conclusion} concludes and points at future work.
